\section{Rigorous local-optimality theorem}\label{sec:local}
\subsection{Validated root and nonsingularity}
Let $c\in\mathbb Q^9$ be the complete decimal centre printed in Appendix~\ref{app:box}, and put $r_0=10^{-20}$. The exact intended box is
\begin{equation}\label{eq:box}
 X=c+[-r_0,r_0]^9.
\end{equation}
Its scale is the coordinate system \eqref{eq:zcoords}, not metres, seconds or a neighborhood on which optimality has been quantified. Outward interval conversion encloses this rational box by a slightly larger representable hull $\widehat X$.

\begin{theorem}[Validated root existence and uniqueness]\label{thm:root}
For the baseline wind \eqref{eq:wind} and the stated A1 force law, $F$ in \eqref{eq:root} has exactly one root $\zeta^*$ in $X$. Every Jacobian enclosed by the interval Jacobian calculation on $\widehat X$ is nonsingular. In particular
\begin{equation}\label{eq:etaexact}\begin{gathered}
 c_6-10^{-20}\leq\eta_0\leq c_6+10^{-20},\\
 c_6=0.152360764301861701918625206455836893973748530913.
\end{gathered}\end{equation}
where $\eta_0=\zeta^*_6$. The three stationary-event coordinates lie strictly in $(0,1)$ and the signed slew branches are $(-1,+1,+1)$.
\end{theorem}
\begin{proof}[Computer-assisted proof]
The force formulas are compositions of polynomials, exponentials, square roots and division with denominators separated from zero on $\widehat X$. They therefore define a $C^2$ map on an open neighborhood of $X$. Interval automatic differentiation supplies an inclusion matrix $\mathbb J(\widehat X)$ containing $DF(\zeta)$ for every $\zeta\in X$. Let $C$ be the fixed matrix of exact finite decimal entries generated by rounding a high-precision inverse at $c$. Its entries are then re-entered as outward intervals; this approximate inverse is not assumed exact.

The verified Krawczyk enclosure is
\begin{equation}\label{eq:krawczyk}
 K(X)=c-CF(c)+(I-C\mathbb J(\widehat X))(X-c).
\end{equation}
The outward computation encloses every component of $K(X)-c$ in $(-1.964\times10^{-34},1.964\times10^{-34})$ and proves
\begin{equation}\label{eq:contraction}
 \sup_{J\in\mathbb J(\widehat X)}\norm{I-CJ}_{\infty}<1.964\times10^{-14}<1.
\end{equation}
The sharper component bounds accompany Table~\ref{tab:rootbox}. These are comparisons of interval endpoints before decimal display rounding. All arithmetic used in the comparisons is outward-rounded at 50 decimal digits.

For any $\zeta\in X$, the integral mean-value formula gives
\[
 \zeta-CF(\zeta)=c-CF(c)+\left(I-C\int_0^1DF(c+t(\zeta-c))dt\right)(\zeta-c)\in K(X).
\]
Thus $\Psi(\zeta)=\zeta-CF(\zeta)$ maps $X$ into its strict interior and, by \eqref{eq:contraction}, is a contraction in the infinity norm. The contraction theorem gives a unique fixed point in $X$. For any enclosed $J$, the matrix $CJ=I-(I-CJ)$ is invertible by its convergent Neumann series. In finite dimensions both $C$ and $J$ are consequently invertible. Hence fixed points of $\Psi$ are precisely zeros of $F$. This proves existence, uniqueness in $X$ and the nonsingularity assertion.

The interval inputs enclose the intended exact $c$ and $r_0$; inclusion is checked against the inner radius, so use of $\widehat X$ does not weaken the conclusion to a different decimal box. The event-time and slew-sign assertions are direct separated interval inequalities. Newton iteration only selected $c$; neither convergence nor its residual supplies any of the existence premises.
\end{proof}

The proof establishes local root uniqueness inside the specified box. It says nothing about other roots outside it. Likewise nonsingularity at fixed baseline wind provides a local implicit branch, not a certified interval of wind values.

\subsection{The active-gradient simplex}
Fix $(s_S,s_L,s_T)$ at their exact root values and write the six functions in \eqref{eq:active} as $g_i(\xi)$ in the five scaled flare coordinates $\xi=(T_f/5,x_D/300,U/75,f_x,f_z)$. Let $G$ have rows $\nabla g_i(\xi^*)^\top$ and define
\begin{equation}\label{eq:simplex}
 A_G=\begin{pmatrix}G^\top\\\mathbf1^\top\end{pmatrix},\qquad
 A_G\lambda=(0,0,0,0,0,1)^\top.
\end{equation}
The proof uses only these fixed-time functions; it does not differentiate an argmin map.

\begin{lemma}[Validated positive simplex]\label{lem:simplex}
The augmented matrix $A_G$ is nonsingular, and its solution $\lambda$ has strictly positive entries. In particular $G$ has rank five, $G^\top\lambda=0$, $\mathbf1^\top\lambda=1$, and zero belongs to the strict interior of the simplex generated by the six gradient vectors.
\end{lemma}
\begin{proof}[Computer-assisted proof]
Apply interval differentiation to the six margins over $\widehat X$, using only the first five derivatives in each row. For a fixed decimal approximate inverse $C_A$, the verified bound is $r_A=\norm{I-C_AA_G}_{\infty}<2.484\times10^{-14}$. The Neumann argument used above proves nonsingularity uniformly over the enclosed matrices. For a chosen centre $\widehat\lambda$ and $b=(0,0,0,0,0,1)^\top$, the residual equation gives
\begin{equation}\label{eq:neumann}
 \norm{\lambda-\widehat\lambda}_{\infty}
 \leq\frac{\norm{C_A(b-A_G\widehat\lambda)}_{\infty}}{1-r_A}.
\end{equation}
The implementation encloses both centre and residual with intervals, so no unvalidated solve is used. The resulting outward summaries are in Table~\ref{tab:weights}; every lower endpoint is positive. Rank five follows because the condition $\operatorname{rank}(G)<5$ would make its five transposed rows, and hence the six rows of $A_G$, dependent. Nonsingularity of $A_G$ also means affine independence of its six gradient columns. Positive barycentric coordinates summing to one put zero in the strict interior of that five-dimensional simplex.
\end{proof}

\begin{table}[htbp]\centering\caption{Outward decimal enclosures of the simplex weights. The source stores longer endpoint strings; these deliberately widened intervals retain strict positivity.}\label{tab:weights}
\begin{tabular}{lll}\toprule
Active margin & Lower endpoint & Upper endpoint\\\midrule
Stopping & $0.8560644515753778$ & $0.8560644515753782$\\
Gate slew & $0.0009658211277946$ & $0.0009658211277950$\\
Touchdown slew & $0.0045094051150227$ & $0.0045094051150231$\\
Interior slew & $0.0087917210360731$ & $0.0087917210360735$\\
Upper load & $0.0514426236725329$ & $0.0514426236725333$\\
Lower thrust & $0.0782259774731976$ & $0.0782259774731981$\\\bottomrule
\end{tabular}\end{table}

\Needspace{12\baselineskip}
\begin{theorem}[Strict local primary optimality in five variables]\label{thm:strict}
Let $p^*$ be the flare coordinates of $\zeta^*$. Define the flare relaxation as the minimum of all continuous flare operating margins and the terminal margins, omitting the approach and hard geometric inequalities. Then $p^*$ is a strict local maximum of this primary objective, with value $\eta_0$, once whole-time feasibility in Theorem~\ref{thm:feasible} is established. There exist $\epsilon,\alpha>0$ such that every nonzero scaled perturbation $h$ with $\norm h<\epsilon$ satisfies
\begin{equation}\label{eq:sharp}
 \eta_f(\xi^*+h)\leq\eta_0-\alpha\norm h<\eta_0.
\end{equation}
No numerical value of $\epsilon$ or $\alpha$ is claimed.
\end{theorem}
\begin{proof}
For nonzero $d$, suppose $\nabla g_i(\xi^*)\cdot d\geq0$ for all $i$. Multiplying by the strictly positive weights and summing gives zero. Each term must therefore vanish. Since $G$ has rank five, $d=0$, a contradiction. Thus every unit vector has a strictly negative minimum inner product.

The function $h_0(d)=\min_i\nabla g_i(\xi^*)\cdot d$ is continuous on the unit sphere. By compactness its maximum is a strictly negative number, say $-\beta$. Hence
\[
 \min_i\nabla g_i(\xi^*)\cdot h\leq-\beta\norm h
 \quad\text{for all }h.
\]
Choose a closed ball around $\xi^*$ in the regular signed-branch domain. The six Hessians are continuous and bounded there, so a common $M<\infty$ satisfies
\[
 |g_i(\xi^*+h)-g_i(\xi^*)-\nabla g_i(\xi^*)\cdot h|
 \leq M\norm h^2.
\]
All six root values equal $\eta_0$. Taking their minimum gives
\[
 \min_i g_i(\xi^*+h)\leq\eta_0-\beta\norm h+M\norm h^2.
\]
For $\norm h<\beta/(2M)$ (with the evident interpretation if $M=0$), this is at most $\eta_0-\beta\norm h/2$. Every full flare minimum is bounded above by these six fixed-time or terminal values. Theorem~\ref{thm:feasible} establishes $\eta_f(\xi^*)=\eta_0$, completing \eqref{eq:sharp} with $\alpha=\beta/2$ and a sufficiently small $\epsilon$.
\end{proof}

This is a first-order sufficient max-min argument. Positive multipliers alone would not suffice without affine independence; necessary KKT conditions alone would not supply the uniform directional decrease. The fixed active-time construction also makes the proof independent of whether other local time extrema move smoothly under perturbation.
\fig{certificate}{Logical structure of the certificate. The diagram depicts implications, not a projection of a five-dimensional simplex. Root isolation, active-gradient sufficiency and continuous-time feasibility supply different premises; only their combination supports the local-optimality conclusion.}{fig:certificate}
